\documentclass{article}
\usepackage[french]{babel}
\usepackage[T1]{fontenc}
\usepackage[utf8]{inputenc}
\usepackage{graphicx} % Required for inserting images
\usepackage{mathtools}
\usepackage{listingsutf8}
\usepackage{amsfonts}
\lstset{%
    inputencoding=utf8,
    literate=
        {é}{{\'e}}{1}%
        {è}{{\`e}}{1}%
        {à}{{\`a}}{1}%
        {â}{{\^a}}{1}%
        {ç}{{\c{c}}}{1}%
        {œ}{{\oe}}{1}%
        {ù}{{\`u}}{1}%
        {É}{{\'E}}{1}%
        {È}{{\`E}}{1}%
        {À}{{\`A}}{1}%
        {Ç}{{\c{C}}}{1}%
        {Œ}{{\OE}}{1}%
        {Ê}{{\^E}}{1}%
        {ê}{{\^e}}{1}%
        {î}{{\^i}}{1}%
        {ï}{{\"i}}{1}%
        {ô}{{\^o}}{1}%
        {û}{{\^u}}{1}%
}
\usepackage{fullpage}
\usepackage[margin=2.1cm]{geometry}
\usepackage{setspace}
\usepackage{todonotes}
\usepackage{amsthm}

\newtheoremstyle{exostyle}% style name
{10pt}% above space
{10pt}% below space
{}% body font
{}% indent amount
{\scshape\bfseries\large}% head font
{\hfill\vspace{5pt}\newline}% post head punctuation
{0pt}% Space after theorem head
{\hfill\thmname{#1}\thmnumber{ #2} -- \thmnote{ #3}}% head spec

\newtheoremstyle{partiestyle}% style name
{1em}% above space
{1em}% below space
{}% body font
{}% indent amount
{\bfseries}% head font
{\vspace{.5em}\newline}% post head punctuation
{0em}% Space after theorem head
{\thmnumber{#2} \thmnote{ #3}}% head spec

\newtheoremstyle{questionstyle}% style name
{.5em}% above space
{.5em}% below space
{}% body font
{}% indent amount
{\bfseries}% head font
{}% post head punctuation
{0em}% Space after theorem head
{Question \thmnumber{#2 }}% head spec

\theoremstyle{exostyle}
\newtheorem{exo}{Exercice}

\theoremstyle{partiestyle}
\newtheorem{partie}{}[exo]

\theoremstyle{questionstyle}
\newtheorem{question}{Question}[exo]
\newtheorem{questionpartie}{Question}[partie]

\title{Devoir surveillé : Algorithmes}
\author{L1 MPCI}
\date{2 mars 2026 - Durée: 2h}

\begin{document}

\maketitle

\begin{center}
{\em\bf Lorsque l'on vous demande d'écrire de décrire ou de donner un algorithme cela signifiera toujours en donner un pseudo-code, justifier de son \underline{exactitude} et de sa \underline{complexité}}

~\\

{\em On rappelle qu'aucun document, ni équipement électrique ou électronique n'est autorisé.

 {\bf Cependant} l'usage d'un marteau-piqueur à percussion sera toléré s'il est utilisé avec un silencieux.}
\end{center}


\vspace*{1cm}

\paragraph{Les exercices de cet examen} :
\begin{itemize}
\item sont au nombre de 3;
\item ont tous le même but, savoir si $b$ est une sous-chaîne de $a$;
\item sont indépendants (à part le problème à résoudre qui est le même);
\item sont de difficulté croissante;
\item sont chacuns constitués de 3 parties de difficultés croissantes.
\end{itemize}

\paragraph{Élément de notation :}
\begin{itemize}
\item  Le nombre de points par partie est capé. {\bf Il est impossible} d'avoir plus de 10/20 (respectivement 16/20) en ne répondant qu'aux questions des premières parties (\textit{resp.} premières et secondes parties) de chaque exercice. {\bf Il n'est cependant pas nécessaire} de répondre à toutes les questions pour avoir 10/20 (\textit{resp.} 16/20);
\item L'examen est {\bf long} et ce qui semble simple pour l'examinateur ne l'est pas forcément pour l'étudiant et réciproquement. Il pourra être utile de changer d'exercice plutôt que de rester bloqué sur une question, quitte à y revenir plus tard.
\end{itemize}


\paragraph{}{\large \sc Rendez \underline{des copies séparées pour chaque exercice}, ceci vous permettra de d'y revenir au cours de l'examen sans perdre le correcteur.}

\vspace*{1cm}

On définit le problème algorithmique suivant :

\paragraph{Problème :}
\begin{itemize}
    \item {\bf Nom :} sous-chaîne
\item {\bf Entrées :}
    \begin{itemize}
        \item \texttt{a} une chaîne de caractères de longueur $n$
        \item \texttt{b} une chaîne de caractères de longueur $m \leq n$
    \end{itemize}
\item {\bf Sortie :} l'indice d'un commencement de $b$ dans $a$ si $b$ est une sous-chaîne de $a$ et $-1$ sinon
\end{itemize}

\paragraph{Avec la définition suivante :} Soient $a$ et $b$ deux chaines de caractères de longueurs $n$ et $m <n$ respectivement. La chaîne $b$ est une \textbf{\textit{sous-chaîne de $a$ commençant en $i$}} si $b = a[i:i+m]$ (c'est à dire que pour $0 \leq i < n-m$ : $b[k] = a[i+k]$ pour tout $0\leq k<m$).


\paragraph{}Par exemple :

\begin{itemize}
\item $corne$, $douille$ et $oui$ sont des sous-chaînes de $cornegidouille$, mais pas $rouille$ ni $ubu$,
\item $a$, $na$ et $an$ sont des sous-chaînes de $nananère$. Elles apparaissent à plusieurs endroits.
\end{itemize}

\clearpage

\begin{exo}[Algorithme Naïf]
    \begin{partie}[Algorithme python]
        \label{algo-python-1}
        \begin{lstlisting}
                def est_sous_chaîne(a, b):
                    for i in range(len(a) - len(b) + 1):
                        trouvé = True
                        j = 0
                        while j < len(b):
                            if b[j] != a[i + j]:
                                trouvé = False
                                j = len(b)
                            else:
                                j += 1
                        if trouvé:
                            return True
                    return False
        \end{lstlisting}
        \begin{questionpartie}
            \label{algo-pseudo-1}
            Donnez le pseudo-code associé à l'algorithme python ci-dessus.
        \end{questionpartie}
        \begin{questionpartie}
            Donnez la complexité minimum et maximum de cet algorithme.
        \end{questionpartie}
        \begin{questionpartie}
            Donnez la signature d'un algorithme résolvant le problème "sous-chaîne".
        \end{questionpartie}
        \begin{questionpartie}
            \label{complexite-naif}
            Modifiez le pseudo-code de la question~\ref{algo-pseudo-1} pour qu'il résolve avec les mêmes complexités le problème "sous-chaîne" (il faudra bien sur démontrer son exactitude).
        \end{questionpartie}
    \end{partie}
    \begin{partie}[Bornes]
        \begin{questionpartie}
            Donnez (et justifiez) une borne minimum au problème "sous-chaîne".
        \end{questionpartie}
        \begin{questionpartie}
            Déduire de la question~\ref{complexite-naif} un encadrement de la complexité du problème "sous-chaîne".
        \end{questionpartie}
    \end{partie}
    \begin{partie}[Complexité en moyenne]
        \begin{questionpartie}
            En supposant l'équiprobabilité des caractères pour les différentes chaînes, montrez que $\Pr [a[i] = b[j]]$ est une constante quelques soient les indices $i$ et $j$
        \end{questionpartie}
        \begin{questionpartie}
            En déduire la complexité en moyenne de l'algorithme de la question~\ref{complexite-naif}\footnote{Vous pourrez utiliser le fait que $\sum_{0\leq i\leq n}(ip^i) \leq \frac{p}{(1-p)^2}$ pour tout $n\geq 0$ si $p<1$}
        \end{questionpartie}

    \end{partie}

\end{exo}
\begin{exo}[Boyer-Moore-Horspool]
    L'algorithme de Boyer-Moore-Horspool  permet de résoudre le problème "sous-chaîne". Il fonctionne de la façon suivante :

    \begin{enumerate}
        \item on commence l'algorithme avec $i= 0$
        \item pour un indice $i$ donné, on fait décroitre un un indice $j$ de $m-1$ à $0$.
        \begin{itemize}
            \item si $a[i + j] = b[j]$ pour tout $j$, l'algorithme s'arrête et rend $i$
            \item sinon, lors de la première non correspondance ($a[i + j] \neq b_j$ et $a[i + j'] = b[j']$ pour $j < j' < m$), l'indice $i$ est {\bf décalé} :
            \begin{itemize}
                \item de $m$ cases si $a[i +m-1]$ n'est pas un caractère de $b[:-1]$ (la chaîne $b$ privée du dernier caractère).
                \item de $m-k-1$ cases où $k$ est l'indice le plus grand dans $b[:-1]$ du caractère $a[i + m -1]$.
            \end{itemize}
        \end{itemize}
        \item si $i> n-m$, l'algorithme s'arrête et rend $-1$ sinon on retourne à l'étape 2
    \end{enumerate}
    \begin{partie}[Principe]
    \begin{questionpartie}
            Démontrez que ce fonctionnement est bien une façon valide de chercher si $b$ est une sous-chaîne de $a$.
    \end{questionpartie}
    \begin{questionpartie}
            En quoi cette amélioration permet d'aller plus vite que l'algorithme basé sur le code de la question~\ref{algo-python-1} ?
    \end{questionpartie}
    \end{partie}
    \begin{partie}[Décalages]
        On suppose que les $L$ caractères possibles  (4 caractères possibles si l'on compare des brins d'ADN, 20 si ce sont des protéines et $2^{28}$ si ce sont des chaînes Unicode)  pour les chaînes $a$ et $b$ sont ordonnés et que l'on possède une fonction de signature \texttt{ord(c: caractère) $\rightarrow$ entier} qui à tout caractère rend son ordre ($0$ pour le plus petit caractère et $L-1$ pour le plus grand) en $\mathcal{O}(1)$ opérations.
        \begin{questionpartie}
            Écrivez une fonction de signature \texttt{décalage(b: chaîne) $\rightarrow$ [entier]} qui rend un tableau $T$ de taille $L$ tel que $T[ord(c)]$ soit égal au décalage à effectuer pour le caractère $c$ lors de la recherche de $b$ dans une chaîne $a$.
        \end{questionpartie}
        \begin{questionpartie}
            Quelle est la complexité de votre fonction \texttt{décalage(b: chaîne) $\rightarrow$ [entier]} ?
        \end{questionpartie}
    \end{partie}
    \begin{partie}[Recherche]
        \begin{questionpartie}
            Écrivez le pseudo-code de l'algorithme de Boyer-Moore-Horspool.
        \end{questionpartie}
        \begin{questionpartie}
            Quelle est la complexité maximale et minimale de l'algorithme de Boyer-Moore-Horspool ?
        \end{questionpartie}
        \begin{questionpartie}
            Quelle est la complexité spatiale de l'algorithme de Boyer-Moore-Horspool ?
        \end{questionpartie}
         \begin{questionpartie}
            En déduire une borne maximum au problème "sous-chaîne".
        \end{questionpartie}
        \begin{questionpartie}
                Quand utiliseriez vous cet algorithme par rapport à l'algorithme naïf ?
        \end{questionpartie}


    \end{partie}
\end{exo}
\begin{exo}[Knuth-Morris-Pratt]
    Cet algorithme est le plus complexe des trois mais, petit {\it spoil}, il est optimal. Tout comme l'algorithme de Boyer-Moore-Horspool il repose sur un tableau de décalage particulier permettant d'accélérer la recherche.
    
        \begin{tabbing}
            ccc\=cccc\=cccc\=cccc\=cccc\=cccc\=\kill
            \textbf{algorithme} \textsc{KMP}($a$: chaîne, $b$: chaîne) $\rightarrow$ entier{\bf :}\\
            \>\vline$\,$ $T \leftarrow$\texttt{décalage}($b$)\\
            \>\vline$\,$ $(i, j \coloneqq \text{entier}) \leftarrow 0, 0$\\
            \>\vline$\,$\textbf{tant que} $i + j < a.\text{longueur}$ {\bf :}\\
            \>\vline$\,$\>\vline$\,$ \textbf{si} $a[i + j] == b[j]${\bf :}\\
            \>\vline$\,$\>\vline$\,$\>\vline$\,$ $j \leftarrow j +1$\\
            \>\vline$\,$\>\vline$\,$ \>\vline$\,$ \textbf{si} $j \geq b.\text{longueur}$ {\bf :}\\
            \>\vline$\,$\>\vline$\,$\>\vline$\,$\>\vline$\,$ \textbf{rendre} $i$\\
            \>\vline$\,$\>\vline$\,$ \textbf{sinon}{\bf :}\\
            \>\vline$\,$\>\vline$\,$ \>\vline$\,$ \textbf{si} $j == 0$ {\bf :}\\
            \>\vline$\,$\>\vline$\,$ \>\vline$\,$\>\vline$\,$ $i \leftarrow i +1$\\
            \>\vline$\,$\>\vline$\,$ \>\vline$\,$ \textbf{sinon}{\bf :}\\
            \>\vline$\,$\>\vline$\,$\>\vline$\,$\>\vline$\,$ $i \leftarrow i + j - T[j]$\\
            \>\vline$\,$\>\vline$\,$\>\vline$\,$\>\vline$\,$ $j \leftarrow T[j]$\\
            \>\vline$\,$ \textbf{rendre} $-1$\\

        \end{tabbing}
         La fonction de décalage $\texttt{décalage}(b: \text{\texttt{chaîne}})\rightarrow [\text{\texttt{entier}}]$ rend un tableau $T$ de même longueur $m$ que $b$ tel que :~\\
            \begin{itemize}
                \item $T$ est une liste d'entier de longueur $m$
                \item $T[0] = T[1] = 0$
                \item pour tout $1 \leq j < m$, $T[j]$ est le plus grand entier $k < j$ tel que $b[:k] = b[j-k:j]$
            \end{itemize}~\\
Par exemple on a $[0, 0, 0, 1, 2, 3, 0] = \text{décalage}("ABABACA")$ :~\\
            \begin{itemize}
                \item $T[2] = 0$ car $b[2-1:2] = "B" \neq b[:1] = "A"$
                \item $T[3] = 1$ car $b[3-1:3] ="A" = b[:1] = "A"$ et $b[3-2:3] ="BA" \neq b[:2] = "AB"$
                \item $T[4] = 2$ car $b[4-2:4] ="AB" = b[:2]$ et $b[4-3:4] ="BAB" \neq b[:3] = "ABA"$
                \item $T[5] = 3$ car $b[5-3:5] ="ABA" = b[:3]$ et $b[5-4:5] ="BABA" \neq b[:4] = "ABAB"$
                \item $T[6] = 0$ car $b[6-1] ="C"$ qui n'est pas dans $b[:5] = "A"$
            \end{itemize}
    \begin{partie}[Complexité Algorithme]
        \begin{questionpartie}
            Montrez qu'à chaque itération soit $i$ soit $i+j$ augmente strictement.
        \end{questionpartie}
        \begin{questionpartie}
            En déduire que la complexité de l'algorithme de Knuth-Morris-Pratt est en $\mathcal{O}(n + f(m))$ avec :
            \begin{itemize}
                \item $n$ et $m$ les longueurs des chaînes $a$ et $b$,
                \item $f(m)$ la complexité de la fonction $\text{décalage}(b)$
            \end{itemize}
        \end{questionpartie}
    \end{partie}
    
    \stepcounter{question}
    \begin{question}
        Montrez l'algorithme de Knuth-Morris-Pratt permet de résoudre le problème "sous-chaîne".
    \end{question}
    \stepcounter{partie}

    \begin{partie}[Décalages]
        La fonction de décalage est extrêmement astucieuse~:\\
        \begin{tabbing}
            ccc\=cccc\=cccc\=cccc\=cccc\=cccc\=\kill
            \textbf{fonction} \textsc{décalage}($b$: chaîne) $\rightarrow$ [entier]{\bf :}\\
            \>\vline$\,$ $(T \coloneqq [\text{entier}]) \leftarrow [\text{entier}]\{\text{longueur:} b.\text{longueur}\}$\\
            \>\vline$\,$ $T[0], T[1] \leftarrow 0, 0$\\
            \>\vline$\,$ $(i, j, k \coloneqq \text{entier}) \leftarrow 2, 2, 0$\\
            \>\vline$\,$ $(c \coloneqq \text{caractère}) \leftarrow b[j-1]$\\
            \>\vline$\,$ \textbf{tant que} $i < T.\text{longueur}$ {\bf :}\\
            \>\vline$\,$\>\vline$\,$ $k \leftarrow T[j-1]$\\
            \>\vline$\,$\>\vline$\,$ \textbf{si } $c == b[k]$ {\bf :}\\
            \>\vline$\,$\>\vline$\,$\>\vline$\,$ $T[i] \leftarrow k + 1$\\
            \>\vline$\,$\>\vline$\,$\>\vline$\,$ $i \leftarrow i + 1$\\
            \>\vline$\,$\>\vline$\,$\>\vline$\,$ $j \leftarrow i$\\
            \>\vline$\,$\>\vline$\,$\>\vline$\,$ $c \leftarrow b[j-1]$\\
            \>\vline$\,$\>\vline$\,$\textbf{sinon} {\bf :} \\
            \>\vline$\,$\>\vline$\,$\>\vline$\,$ \textbf{si } $k \leq 1$ {\bf :}\\
            \>\vline$\,$\>\vline$\,$\>\vline$\,$\>\vline$\,$ $T[i] \leftarrow 0$\\
            \>\vline$\,$\>\vline$\,$\>\vline$\,$\>\vline$\,$ $i \leftarrow i + 1$\\
            \>\vline$\,$\>\vline$\,$\>\vline$\,$\>\vline$\,$ $j \leftarrow i$\\
            \>\vline$\,$\>\vline$\,$\>\vline$\,$\>\vline$\,$ $c \leftarrow b[j-1]$\\
            \>\vline$\,$\>\vline$\,$\>\vline$\,$ \textbf{sinon} {\bf :}\\
            \>\vline$\,$\>\vline$\,$\>\vline$\,$\>\vline$\,$ $j \leftarrow k + 1$\\
            \>\vline$\,$ \textbf{rendre} $T$\\
        \end{tabbing}
        \begin{questionpartie}
            Montrez qu'à chaque itération soit :
            \begin{itemize}
                \item $k$ reste constant à 0
                \item $k$ augmente de 1
                \item $k$ diminue strictement
            \end{itemize}
        \end{questionpartie}
        \begin{questionpartie}
            En déduire que la complexité de la fonction \texttt{décalage}(b) est en $\mathcal{O}(m)$ avec $m$ la taille de la chaîne passée en paramètre de la fonction.
        \end{questionpartie}
        \begin{questionpartie}
            Montrez que la fonction décalage correspond bien à celle demandée pour faire fonctionner l'algorithme de Knuth-Morris-Pratt.
        \end{questionpartie}
    \end{partie}
\end{exo}
\begin{exo}[Exercice secret final et bonus]
    Déduire de tout ce qui précède la complexité de problème "sous-chaîne".
\end{exo}

\end{document}